\section{Borrowed wires and the cost of an error}\label{sec:catalysts}

We now use a maximal $k$-flat to route a permutation. Each payload must
reach its prescribed output. Alongside the payloads we supply pairs of
binary ports at opposite ends of the coordinate axes and require the
router to return an identity wire on every pair. These extra wires are
the \emph{catalysts}. Their placements have two roles. In the construction
of Section~\ref{sec:protocol}, keeping all $k$ directions at full price
constrains the local wire distances available to the router. In
Section~\ref{sec:width}, a term that saves sufficiently much on the
catalysts may be discarded, and the surviving terms leave little budget
for transporting payloads. We define the permutation tensor and make
precise how much error this argument permits.

\paragraph{Bit networks.}
A bit network has a binary index at each port. A wire equates two
indices and sums over their common bit. The entry at an assignment of
the open bits is the sum, over internal bits, of the product of the local
tensor entries. A trivalent network has at most three ports per local
tensor. Size counts tensors and wires with multiplicity, including a wire
for each open port. Equivalently, an open port can be a terminal wire end,
counting both ends when both are open. Joined ends cease to be open and may
be joined without subdividing the edge. A closed wire with no vertex
contributes the field element $2$.

Any fixed arity reduces to three at constant cost by expanding the local
table in elementary tensors and testing the selected entry. Bit copying
and Boolean gates have trivalent implementations, with the cost of
fan-out included. In particular, a deterministic verification computes
the indicator tensor of its relation over any field if each accepting
input has exactly one witness, including its padding bits. Unless
specified otherwise, we bound network size without a uniformity
requirement.

\begin{definition}[Tagged permutation]\label{def:tag}
For integers $h,w\ge1$, put $n=2^h$. The open data of
$\operatorname{Tag}(h,w)$ are an ordered list $(p_a)_{a=0}^{n-1}$ of
$w$-bit payloads and a list $((a_i,q_i))_{i=0}^{n-1}$ of records with
$h$-bit address and $w$-bit payload. The tensor entry is one exactly when
the addresses $a_i$ are a permutation of $0,\ldots,n-1$ and
$q_i=p_{a_i}$ for every $i$. All other entries are zero.
\end{definition}

\subsection{Approximation on an ordered Euclidean frame}

Let $e_1,\ldots,e_k$ be the coordinate vectors. A tensor $X$ has data ports
at the origin and $b_j$ pairs of catalyst ports, with one port per pair at
each of $-e_j$ and $+e_j$. Write $I_X$ for the product of identity wires
on these pairs and $h_X=2\sum_j b_j$ for its price. We say that $X$
\emph{approximates} a data tensor $f$ with \emph{gap} $D$ if
\begin{equation}\label{eq:catalyst-gap}
 H(X-f\otimes I_X)\le h_X-D
\end{equation}
for every isometric placement of this ordered frame induced by a spectral
$k$-flat over $F$, in any spectral space. The estimate consequently holds
at translated and rotated frames and, with both sides multiplied by the
dilation factor, on dilated frames. The latter statement follows by first
rescaling the entire spectral space. Thus $D$ measures the saving in
price of the error relative to the exact return of the borrowed wires.
If $X=f\otimes I_X$, every finite gap is admissible. All open catalyst
bits contribute to network size.

Exposed data ports in \eqref{eq:catalyst-gap} remain available for arbitrary
local contractions. For designated \emph{pins} we sometimes require the
estimate only after fixing all their bit values. Exposing $p$ such pins costs
at most $2kp$ in gap. In fact the full tensor is the sum of the $2^p$
conditional tensors tensored with their local pin indicators. The star
joining $0$ to the $2k$ catalyst sites bounds the summation scale by $2k$.
Fixing an exposed bit is a restriction and cannot increase the error price.

The following estimate accounts for the loss of gap when several
approximations are used in the same contraction.
\begin{lemma}[Substitution]\label{lem:substitution}
Suppose a contraction approximates $g$ with gap $D$ using $m$ exact data
tensors $f_i$, counted by occurrence. Replace occurrence $i$ by
$X_i$ of gap $D_i\ge0$ relative to $f_i$, keeping all new catalyst pairs
open. The resulting gap is at least
$\min(D,D_1,\ldots,D_m)-2km$. Pointwise estimates can be used on any pins
whose values stay fixed throughout the contraction.
\end{lemma}
\begin{proof}
Expand each replacement as $f_i\otimes I_{X_i}$ plus its error. In the term
with no replacement errors, subtract the target. The original error bound
applies with the prices of the new identity wires added. For a nonempty set of
replacement errors, carry out the old contraction at the origin, placing the
new catalyst ports in their prescribed frames. At every erroneous occurrence
save $D_i$ from its through price. Other new pairs cost only their through
prices. Move the old exterior catalyst bits from the origin to their
prescribed sites. This costs one per bit, in total exactly the budget for the
old through wires. Thus every term is below the combined through budget by at
least the indicated minimum. There are at most $2^m$ terms, and
\eqref{eq:price-ordinary-sum} with the star bound proves the result.
\end{proof}
